\section{Discussion}
This project could be seen as an introduction to the VESC project for someone who don't know about it from beforehand, the challenges the new users face during setup, as well as a demand for clear expectations concerning documentation on the subject. The project the MAD is leading should probably not be a fork of the project, as the project is still in development.

As a final note, this proved to be a project which could easily be developed into several different projects in different fields. Some projects could be continued later on as a different PIR subject, other could be proposed to later years in different specialisations like TLS-SEC, ESPE. Our thoughts on the following projects that could be 

The fabrication line for electronics is globalised. This is okay in a stable world, but it could be a problem in a world full of instability, be it war, blockages, or tariffs. The idea of opening a specialisation in cooperation with AIME came up as an idea.

For TLS-SEC the subject could be the design for a fitting mechanism to restrict certain privileges to certified personnel that could be used in the C programming language. Later down the line we could also see the possibility to analyse the Bluetooth frames in order to manipulate them in order to change important parameters.

The continuation on the PCB could be a subject fitting an ESPE specialisation.

The proposition of and supply of a VESC system to play with and troubleshoot could be a good rule of thumb, which allows for a quicker start and gives among other things an idea of the budget and the supply line used by a entity in the sector. Proposing a visit could also be one way to familiarise students with the association.

What should be a clear conclusion from our test with the jammer is that a controller based on Bluetooth alone should be avoided when possible and practical. Examples where this could be relevant include electric skateboards, as cables could impose a tripping hazard. There, an encapsulation of an encrypted control frame could be an thought. 



